\setlength{\headheight}{15pt}   % NOTE: not sure if this is needed.
\chapter*{Abstract}
\setcounter{page}{1}  % <-- reset so we do not count title
\addcontentsline{toc}{chapter}{Abstract} % Add the abstract to page of contents

% \thispagestyle{fancy}
% A summary of no more than 300 words;
% https://wordcounter.net/

% Social media has become an integral part of daily life for millions worldwide, serving as a platform where people exchange opinions, express concerns, and influence one another. From natural disasters to political debates, social media platforms like X (formerly Twitter) and Reddit have become primary sources of news and public discourse. This activity results in the generation of vast amounts of data each day which, due to their sheer volume and largely unstructured nature, researchers and analysts face significant challenges in effectively studying and analysing it.

% To analyse this complex data source, researchers frequently utilise \ac{NLP} tools and techniques. However, social media data poses unique challenges distinct from more formal or structured text. The landscape of social media is constantly evolving, with new slang, trending topics, emojis, and brevity of posts creating additional hurdles that traditional \ac{NLP} methods often struggle to address.

% Despite the impact of models such as BERT and RoBERTa, and more recently \acp{LLM} like GPT and LLaMA, social media data continues to present significant difficulties due to its distinctive characteristics. While large \acp{LLM} have demonstrated broad and impressive capabilities across a range of tasks they are not always the most suitable choice. Architectural differences between these decoder-based models and smaller encoder-based models can make the latter more efficient and effective for certain domain-specific applications.

% In this work, we present a collection of high-quality, domain-specific datasets sourced from social media, designed to support both the training and evaluation of targeted \ac{NLP} models. Building on these resources, we develop small encoder-based models tailored to specific tasks and demonstrate that their performance can match, and in some cases exceed, that of larger decoder-based \acp{LLM} on domain-specific benchmarks. These tools are applied to case studies involving politically charged discourse, hate speech detection, and misinformation analysis. Collectively, the datasets and models presented in this work highlight that specialised and efficient approaches can achieve competitive performance while supporting scalable and sustainable social media research.

% NOTE: above was 325 words so we rephrase a bit
Social media has become central to daily life for millions worldwide, serving as a space where people exchange opinions, express concerns, and influence one another. From natural disasters to political debates, platforms such as X (formerly Twitter) and Reddit are now primary sources of news and public discourse. This activity generates vast amounts of largely unstructured data, creating major challenges for researchers seeking to study and analyse it effectively.

\ac{NLP} offers tools to handle this complexity, but social media poses unique difficulties compared to more formal or structured text. Rapidly evolving slang, trending topics, emojis, and the brevity of posts create hurdles that traditional \ac{NLP} methods often struggle to overcome.

Despite the impact of models such as BERT and RoBERTa, and more recently \acp{LLM} like GPT and LLaMA, social media data's distinct characteristics continue to pose difficulties. While \acp{LLM} show impressive, broad capabilities, architectural differences mean that smaller, encoder-based models can be more efficient and effective for certain domain-specific applications.

In this work, we present a collection of high-quality, domain-specific datasets sourced from social media to support the training and evaluation of targeted \ac{NLP} models. Building on these resources, we develop small encoder-based models tailored to specific tasks, demonstrating that their performance can match, and in some cases exceed, that of larger decoder-based \acp{LLM} on domain-specific benchmarks. These tools are applied to case studies involving politically charged discourse, hate speech detection, and misinformation analysis. Collectively, this research highlights that specialised and efficient approaches can achieve competitive performance while supporting scalable and sustainable social media research.